
\documentclass[11pt]{article}
\usepackage{times}
\usepackage{url}
\usepackage{latexsym}
\usepackage[utf8]{inputenc}  %% to get umlauts right
\usepackage{graphicx}
\usepackage{covington}
\usepackage{natbib}

\usepackage{breakcites}
 
\setlength{\topmargin}{-1cm}
\setlength{\textheight}{23cm}
\setlength{\oddsidemargin}{0.0cm}
\setlength{\textwidth}{16cm}


\pagestyle{myheadings}    % Go for customized headings
\markboth{Test}{Names: Project title \hspace{2cm} \today \hfill }


\begin{document}


\begin{titlepage}
\includegraphics[height=15mm]{uzh_logo_d_pos}\\


\begin{center}

{\sffamily {\Large
Seminar \\
\textbf{Model Analysis and Interpretability\\ for Natural Language Processing} \\
Spring term 2026 \\

\vspace{2cm}

{\huge Title}\\

\vspace{4cm}

\textbf{Author: Zhaoyan, Fan} \\
	Student Id: xx-xxx-xxx   \\

\vspace{4cm}
}
Course instructor: Rico Sennrich\\
Department of Computational Linguistics

\vfill Submission date: (02.06.2026) 

\vspace{3cm}
}
\end{center}

\end{titlepage}

\newpage

Please use this template for your term paper. You can copy it by selecting "copy" in the list view of your projects. It should be max. 8 pages long including tables and figures but excluding the list of references. The following structure is only a suggestion, feel free to change it if that makes sense for your project.

\section{Introduction}

Your motivation for the topic and your research question(s) should go here. It can include a brief summary of your results and conclusions.

\section{Related Work}

A summary of related research should go here. It can be structured with subsections. Please references with author name and year: \citep{hale2001probabilistic} or \citet{hale2001probabilistic}.

\section{Model architectures}
If you've trained/evaluated a model/ different models, describe their architecture in this section.

\section{Experiments}
In a typical paper the Experiments section contains a description of models, data, and evaluation methods used in your experiments.


\section{Results}
Please include result tables and figures here. In of psycholinguistic papers, the Results constitute a separate section. In Computational Linguistics, the Results are usually part of the Experiments section. 

\section{Discussion}
A discussion of your findings and its implications for the field should go here. Importantly, you should also discuss the limitations of your work. 
 
 
\section{Conclusion}
A very short summary of your work goes here: What is the take-home-message that you want the reader to remember? 


\bibliography{custom}
\bibliographystyle{acl_natbib}

\end{document}
